\chapter{Projective connections and three-dimensional geometry}
\label{ch:projective-gravity}

A second-order equation records how two independent solutions change along a curve.
Their ratio removes the choice of a common scale and retains a projective coordinate.
Two such equations, one for each boundary direction, can determine a three-dimensional metric.
The metric exists where the corresponding three one-forms are independent.
This independence condition distinguishes a gravitational geometry from a pair of flat matrix connections.

The quadratic central current supplies the coefficient of the differential equation.
The construction therefore has three substantive steps: evaluate that current, retain its coordinate transformation, and construct a nondegenerate metric.
The constant-coefficient case will also determine a regular extension through a horizon.

\section{The central character and its differential equation}

Let $z$ be a complex coordinate on a disc punctured at $0$.
It measures position along the curve on which the fields are evaluated.
It is independent of the coordinate on the cotangent surface.
Write $\partial_z$ for differentiation with respect to $z$.
For a meromorphic function $\chi(z)$ with a pole of finite order, consider
\[
 \mathcal U(z)=\frac{\chi(z)^2}{4}-\frac{\partial_z\chi(z)}2,
 \qquad \mathscr D_z=\partial_z^2-\mathcal U(z).
\]
The product of two differential operators means composition, with the right factor acting first.
The product rule gives
\begin{equation}\label{eq:pg-miura}
 \mathscr D_z=
 \left(\partial_z-\frac{\chi(z)}2\right)
 \left(\partial_z+\frac{\chi(z)}2\right).
\end{equation}
The derivative term is forced by this order.

This coefficient is an evaluation of the critical quadratic current.
At affine level $-2$, the field $a$ has commuting modes and zero singular product with every field.
The exact quadratic expression from \eqref{eq:cc-miura} is
\[
 S=\tfrac12:a^2:-Ta.
\]
Here $T$ differentiates fields in $z$, and the colons denote the regular product.
Expand $a(z)=\sum_{n\in\mathbb Z}a_nz^{-n-1}$.
The assignment $a_n\mapsto\chi_n$ is a representation of its commuting mode algebra when
$\chi(z)=\sum_n\chi_nz^{-n-1}$.
The finite pole order ensures $\chi_n=0$ for sufficiently large $n$.
Thus each field acts as a Laurent series bounded below on each vector.
Tensor this one-dimensional mode representation with the polynomial $\beta\gamma$ representation.
The critical current formulas still act: all commutators involving the substituted $a_n$ vanish, as required.
Products and derivatives of central fields act by products and derivatives of $\chi$.
Consequently
\[
 S(z)\longmapsto\tfrac12\chi(z)^2-\partial_z\chi(z)=2\mathcal U(z).
\]
This is a character on a module; it preserves the positional dependence of the field.

For $\nu\in\mathbb C$, the mode character $a_0=\nu$, $a_n=0$ for $n\ne0$, gives
\begin{equation}\label{eq:pg-euler}
 \chi(z)=\frac\nu z,
 \qquad \mathcal U_\nu(z)=\frac{\nu(\nu+2)}{4z^2},
 \qquad \mathscr D_z\psi=0.
\end{equation}
The scalar $\nu(\nu+2)$ is the quadratic central character of the differential-operator module already constructed.
The shift by $2$ comes from differentiating $\nu/z$ in \eqref{eq:pg-miura}.
Replacing the factors by commuting symbols would instead give $\nu^2$.

Choose a branch of $\log z$ on a simply connected part of the punctured disc.
Substitution of $\psi=z^\alpha=\exp(\alpha\log z)$ into \eqref{eq:pg-euler} gives
\[
 \alpha(\alpha-1)=\frac{\nu(\nu+2)}4.
\]
The roots are $(\nu+2)/2$ and $-\nu/2$.
The Wronskian of functions $\psi_1,\psi_2$ means $\psi_1'\psi_2-\psi_1\psi_2'$.
When $\nu\ne-1$, the corresponding powers are independent because their Wronskian is a nonzero constant.
Its derivative is zero for two solutions of this equation.
At $\nu=-1$, substitution verifies the independent solutions
\[
 z^{1/2},\qquad z^{1/2}\log z.
\]
Thus an integral difference between distinct exponents produces no logarithm in this Euler equation.
The logarithm occurs at the repeated exponent.

Analytic continuation once counterclockwise around $0$ is called monodromy.
For $\nu\ne-1$, its eigenvalues on the displayed powers are
$e^{\pi i\nu}$ and $e^{-\pi i\nu}$.
At $\nu=-1$, it acts on the ordered solution row by
\[
 (\psi_1,\psi_2)\longmapsto
 (\psi_1,\psi_2)
 \begin{pmatrix}-1&-2\pi i\\0&-1\end{pmatrix}.
\]
The reflection $\nu\mapsto-\nu-2$ preserves the equation and exchanges the two distinct exponents.
It retains the differential equation while changing the selected first-order factor.
The invariant $\nu(\nu+2)=(\nu+1)^2-1$ is unchanged when $\nu+1$ changes sign.
At the fixed point $\nu=-1$, the coalescence of the exponents requires the logarithmic solution to retain the full monodromy.

\section{Why the coefficient is a projective connection}

The equation $\psi''=\mathcal U\psi$ can be written as a first-order matrix equation
\[
 \frac{d}{dz}\begin{pmatrix}\psi'\\\psi\end{pmatrix}
 =\begin{pmatrix}0&\mathcal U\\1&0\end{pmatrix}
  \begin{pmatrix}\psi'\\\psi\end{pmatrix}.
\]
If $\mathcal U$ is holomorphic on a disc, prescribed initial values determine a unique holomorphic solution.
Here is a direct existence argument.
For the matrix coefficient $M(z)$, iterate
$Y(z)=I+\int_{z_0}^zM(w)Y(w)\,dw$ along a segment in a smaller disc.
If $\|M\|\le C$ and its radius is $d$, the term with $n$ integrals has norm at most $(Cd)^n/n!$.
The resulting series and its differentiated series converge uniformly on smaller discs.
They give $Y'=MY$ and $Y(z_0)=I$.
Applying the same iterated estimate to a difference proves uniqueness.
The trace of $M$ is zero, so differentiating the two-column determinant gives $\det Y=1$.
Continuation through overlapping discs therefore gives a two-dimensional local solution space with preserved Wronskian.

Let $z=f(t)$ be a holomorphic change of coordinate with $f'(t)\ne0$.
Locally choose a square root of $f'$ and set
\[
 \widetilde\psi(t)=(f'(t))^{-1/2}\psi(f(t)).
\]
The exponent $-1/2$ is forced by the requirement that the transformed equation have no first derivative.
Indeed, for a general multiplier $h(t)$, the coefficient of $\psi'(f(t))$ in
$(h\psi(f))''$ is $2h'f'+hf''$.
It vanishes exactly when $h$ is a constant multiple of $(f')^{-1/2}$.
Such transformation laws define a section of the inverse square root of the cotangent line, also called a differential of weight $-1/2$.
Square roots on several charts require compatible choices of signs.
The projective equation below does not depend on those signs.

Define the Schwarzian derivative by
\[
 \{f,t\}=\frac{f'''}{f'}-\frac32\left(\frac{f''}{f'}\right)^2.
\]
Differentiating $h=(f')^{-1/2}$ gives $h''/h=-\{f,t\}/2$.
Hence
\begin{equation}\label{eq:pg-projective-law}
 \widetilde\psi''-\widetilde{\mathcal U}\widetilde\psi
 =(f')^{3/2}(\psi''-\mathcal U\psi)\circ f,
 \qquad
 \widetilde{\mathcal U}=(f')^2\mathcal U\circ f-\tfrac12\{f,t\}.
\end{equation}
A projective connection is a collection of local coefficients related by this transformation law.
The chain rule for third derivatives gives
\[
 \{f\circ g,s\}=(g')^2\{f,g(s)\}+\{g,s\}.
\]
Thus the law composes on triple overlaps.
The difference of two projective connections transforms by multiplication by $(f')^2$.
It is therefore a quadratic differential, meaning a section locally written $q(z)(dz)^2$.

A projective coordinate is a local function with nonzero derivative, considered up to the fractional linear changes $(aF+b)/(cF+d)$ with $ad-bc\ne0$.

\begin{proposition}\label{prop:pg-projective-ratio}
The ratio of two independent solutions gives local projective coordinates whose Schwarzian is $-2\mathcal U$.
Every holomorphic projective connection admits a local factorization \eqref{eq:pg-miura}.
The choice of factorization is the choice of a nonvanishing solution line.
\end{proposition}
\begin{proof}
Let $F=\psi_1/\psi_2$ on a region where $\psi_2\ne0$.
Writing its nonzero constant Wronskian as $W$ gives
$F'=W/\psi_2^2$ and $F''/F'=-2\psi_2'/\psi_2$.
Substitution in the Schwarzian yields $\{F,z\}=-2\psi_2''/\psi_2=-2\mathcal U$.
A change of solution basis sends $F$ to $(aF+b)/(cF+d)$, with $ad-bc\ne0$.
Its Schwarzian is unchanged, as direct differentiation of this fraction shows.
These fractional linear transformations are the changes of projective coordinate.

Choose a solution $\psi_0$ with $\psi_0(z_0)=1$.
It remains nonzero after shrinking the disc.
Set $\chi=-2\psi_0'/\psi_0$.
Then $\chi^2/4-\chi'/2=\psi_0''/\psi_0=\mathcal U$.
Multiplying $\psi_0$ by a nonzero constant does not change $\chi$.
Conversely, the first-order equation $(\partial_z+\chi/2)\psi_0=0$
has precisely one solution line on a simply connected disc.
This proves the last assertion.
\end{proof}

The first-order coefficient has the transformation law
\[
 \widetilde\chi=f'\chi\circ f+\frac{f''}{f'}.
\]
Inserting it into \eqref{eq:pg-miura} gives \eqref{eq:pg-projective-law}, including the Schwarzian term.
This is the rank-two Miura construction in explicit coordinates.
The solution columns transform by the explicit matrix
\[
 \begin{pmatrix}\widetilde\psi'\\\widetilde\psi\end{pmatrix}
 =\begin{pmatrix}(f')^{1/2}&-f''/(2(f')^{3/2})\\0&(f')^{-1/2}\end{pmatrix}
   \begin{pmatrix}\psi'\circ f\\\psi\circ f\end{pmatrix}.
\]
Its determinant is one, and its upper triangular form preserves the first coordinate line.
These matrices glue the local two-dimensional vector spaces and their connection equations into a rank-two bundle with a flat connection.
A line subbundle chooses a one-dimensional subspace varying holomorphically in each fibre.
A rank-two oper here is such a determinant-one flat bundle with a line subbundle whose derivative maps isomorphically to the quotient tensored with the cotangent line.
In the displayed first-order system, take the first coordinate line.
Its derivative has nonzero second component, so it has this property.
The half-differential transformations give the same local scalar operator description.
A selected flat solution line is the additional Miura datum.
A zero of the selected solution produces a pole in $\chi$.
A global holomorphic factorization therefore requires a nowhere-zero solution line preserved by monodromy.
These requirements can obstruct a factorization on a prescribed domain without additional poles.

For \eqref{eq:pg-euler}, the projective coordinates can be chosen as
\[
 F(z)=z^{\nu+1}\quad(\nu\ne-1),\qquad F(z)=\log z\quad(\nu=-1).
\]
The repeated exponent therefore has nontrivial projective monodromy even though both linear eigenvalues coincide.

\section{Two real flat connections}

A real Lorentzian metric needs both a time direction and a space direction.
Supply two smooth real functions $p(x^+)$ and $q(x^-)$ on real intervals, and a length $\ell>0$.
The functions may be restrictions of \eqref{eq:pg-euler} with real $\nu$ on positive intervals.
They may also be chosen independently.
Introduce a dimensionless radial coordinate $r$.
All three coordinates $(r,x^+,x^-)$ vary independently.

Let $\mathfrak{sl}_2(\mathbb R)$ be the traceless real $2\times2$ matrices, with bracket $[X,Y]=XY-YX$.
Use the basis
\[
 L_0=\frac12\begin{pmatrix}1&0\\0&-1\end{pmatrix},\qquad
 L_1=\begin{pmatrix}0&0\\-1&0\end{pmatrix},\qquad
 L_{-1}=\begin{pmatrix}0&1\\0&0\end{pmatrix}.
\]
Matrix multiplication gives $[L_m,L_n]=(m-n)L_{m+n}$ whenever the displayed indices occur.
A matrix connection is an operator $d+A$ on column-valued functions, where $A$ is a matrix-valued one-form.
Its curvature is $F_A=dA+A\wedge A$.
The exterior product combines matrix multiplication with the alternating product of coordinate differentials.
Thus $(A\wedge A)_{ij}=[A_i,A_j]$ for two different coordinate indices.
The connection is flat when $F_A=0$.

Put $b=\exp(rL_0)$ and define
\begin{align}
 A^+&=b^{-1}(L_1-pL_{-1})b\,dx^++b^{-1}db,
       \label{eq:pg-connections}\\
 A^-&=b(L_{-1}-qL_1)b^{-1}\,dx^-+b\,db^{-1}.\notag
\end{align}
The untransformed one-form $(L_1-p(x^+)L_{-1})dx^+$ has zero exterior derivative and zero exterior square.
The same holds for the minus connection.
Under $A\mapsto g^{-1}Ag+g^{-1}dg$, direct differentiation gives
$F_A\mapsto g^{-1}F_Ag$.
Both connections in \eqref{eq:pg-connections} are therefore flat.

The scalar equations enter through actual horizontal sections.
The equation $(d+(L_1-pL_{-1})dx^+)v=0$ is solved by
$v=(\psi',\psi)^t$ exactly when $\psi''=p\psi$.
For the minus connection, the column $(\psi,-\psi')^t$ is horizontal exactly when $\psi''=q\psi$.
Multiplication by $b^{-1}$ or $b$ gives the corresponding radial horizontal sections.
Thus the original projective equations determine the boundary parallel transport of these specified connections.

\section{The metric and its curvature}

The difference of the two connections defines a matrix-valued one-form
\[
 e=\frac\ell2(A^+-A^-),\qquad
 \omega=\frac12(A^++A^-).
\]
The form $e$ is a triad when the map $e:T_xM\to\mathfrak{sl}_2(\mathbb R)$ is invertible at every point.
It identifies a tangent vector with its components in a local Lorentz frame.
The invariant bilinear form $B(X,Y)=2\operatorname{Tr}(XY)$ has signature $(-,+,+)$.
For example, the basis
\[
 J_0=\tfrac12(L_1+L_{-1}),\qquad
 J_1=\tfrac12(L_{-1}-L_1),\qquad J_2=L_0
\]
has $B(J_a,J_b)=\operatorname{diag}(-1,1,1)_{ab}$.
Consequently an invertible triad defines a Lorentzian metric
$g(V,W)=B(e(V),e(W))$.

Write $s=e^r>0$.
The three columns of $e$ in the ordered basis $(L_0,L_1,L_{-1})$ are
\[
 e_r=\ell L_0,\qquad
 e_+=\frac\ell2(sL_1-ps^{-1}L_{-1}),\qquad
 e_-=\frac\ell2(qs^{-1}L_1-sL_{-1}).
\]
Their determinant and their metric are
\begin{align}
 \det(e_r,e_+,e_-)&=-\frac{\ell^3}{4}(s^2-pq\,s^{-2}),
       \label{eq:pg-triad}\\
 ds_g^2&=\ell^2\bigl(dr^2+p(dx^+)^2+q(dx^-)^2
                  +(s^2+pq\,s^{-2})dx^+dx^-\bigr).
       \label{eq:pg-metric}
\end{align}
Here $ds_g^2$ denotes the metric's quadratic form; the mixed coefficient is twice $g_{+-}$.
Its determinant in these coordinates is
$-\ell^6(s^2-pq\,s^{-2})^2/4$.

For a connection $\nabla$ on tangent vectors, use the curvature convention
\[
 \mathcal R(V,W)Z=\nabla_V\nabla_WZ-\nabla_W\nabla_VZ-\nabla_{[V,W]}Z.
\]
Its Ricci tensor is $\operatorname{Ric}(W,Z)=\operatorname{Tr}(V\mapsto\mathcal R(V,W)Z)$.
Its scalar curvature is the contraction of this tensor with the inverse metric.
Constant sectional curvature $K$ means
$\mathcal R(V,W)Z=K(g(W,Z)V-g(V,Z)W)$.
These definitions apply to Lorentzian metrics as well as positive definite metrics.

\begin{theorem}\label{thm:pg-local-einstein}
On the open set $s^4\ne pq$, formula \eqref{eq:pg-metric} is a smooth Lorentzian metric of constant sectional curvature $-\ell^{-2}$.
Its Ricci tensor is $-2\ell^{-2}g$ and its scalar curvature is $-6\ell^{-2}$.
Every compact pair of boundary intervals admits a nondegenerate collar at sufficiently large $r$.
\end{theorem}
\begin{proof}
Adding and subtracting the flatness equations gives
\begin{equation}\label{eq:pg-frame-equations}
 de+\omega\wedge e+e\wedge\omega=0,
 \qquad F_\omega+\ell^{-2}e\wedge e=0.
\end{equation}
The first equation says that the covariant derivative
\[
 \nabla_VW=e^{-1}\bigl(V(e(W))+[\omega(V),e(W)]\bigr)
\]
has zero torsion: $\nabla_VW-\nabla_WV=[V,W]$.
The trace identity $B([X,Y],Z)+B(Y,[X,Z])=0$ proves compatibility with $g$.
This is the Levi--Civita connection.
For uniqueness, the difference of two metric-compatible connections is skew in its last two lowered indices.
Zero torsion makes it symmetric in the first two.
Cycling these symmetries changes its sign without changing the tensor, so the difference vanishes.

The matrix identity
\[
 [[X,Y],Z]=B(Y,Z)X-B(X,Z)Y
\]
follows by evaluating the three basis matrices and extending trilinearly.
The second equation in \eqref{eq:pg-frame-equations} therefore gives
\[
 \mathcal R(V,W)Z=-\ell^{-2}\bigl(g(W,Z)V-g(V,Z)W\bigr).
\]
This is constant curvature $-\ell^{-2}$, with the stated contractions.
Finally, $pq$ is bounded on compact boundary intervals.
Choose $r_0$ with $e^{4r_0}$ larger than that bound.
Equation \eqref{eq:pg-triad} is then nonzero for $r>r_0$.
\end{proof}

For $p=q=0$, introduce time $t=\ell(x^+-x^-)/2$ and space $\varphi=(x^++x^-)/2$.
The metric becomes
\[
 ds_g^2=\ell^2dr^2+e^{2r}(\ell^2d\varphi^2-dt^2).
\]
This fixes the sign of the time direction in \eqref{eq:pg-metric}.
In particular, its positive $dx^+dx^-$ coefficient is consistent with Lorentzian signature.

\section{The gravitational action and its normalization}

An action determines which variations represent the equations of motion and fixes the coupling of boundary observables.
Choose an oriented compact region in the nondegenerate collar, and restrict variations to its interior.
Write $e=\theta^aJ_a$ and choose the orientation for which
$\mathrm{vol}_e=\theta^0\wedge\theta^1\wedge\theta^2$ is positive.
Repeated frame indices run from $0$ to $2$.
The one-forms $\theta^a$ constitute an orthonormal coframe for $g$.
Let $G>0$ have dimensions of length; in three dimensions it is Newton's constant.
Set
\[
 I_{\mathrm P}[e,\omega]
 =\frac1{4\pi G}\int\operatorname{Tr}
       \left(e\wedge F_\omega+\frac1{3\ell^2}e\wedge e\wedge e\right).
\]
The subscript denotes the first-order, or Palatini, action in which $e$ and $\omega$ vary independently.

\begin{proposition}\label{prop:pg-action}
Interior variations of $I_{\mathrm P}$ give exactly \eqref{eq:pg-frame-equations}.
After eliminating the torsion-free connection, the action is
\[
 I_{\mathrm{EH}}[g]=\frac1{16\pi G}\int
          (\mathrm{Scal}(g)+2\ell^{-2})\,\mathrm{vol}_g.
\]
Define $\mathrm{CS}(A)=\operatorname{Tr}(A\wedge dA+\tfrac23A\wedge A\wedge A)$.
For $k=\ell/(4G)$ and $A^\pm=\omega\pm e/\ell$, the exact identity of actions is
\begin{equation}\label{eq:pg-action-comparison}
 \frac{k}{4\pi}\int\bigl(\mathrm{CS}(A^+)-\mathrm{CS}(A^-)\bigr)
 =I_{\mathrm P}-\frac1{8\pi G}\int_{\partial M}\operatorname{Tr}(\omega\wedge e).
\end{equation}
\end{proposition}
\begin{proof}
The variation of $F_\omega$ is $D_\omega\delta\omega$.
The identity
$d\operatorname{Tr}(e\wedge\delta\omega)
=\operatorname{Tr}(D_\omega e\wedge\delta\omega-e\wedge D_\omega\delta\omega)$
integrates the connection variation by parts.
Cyclicity of the trace makes the cubic variation three times
$\operatorname{Tr}(\delta e\wedge e\wedge e)$.
The interior variation is therefore
\[
 4\pi G\,\delta I_{\mathrm P}
 =\int\operatorname{Tr}\bigl(
   \delta e\wedge(F_\omega+\ell^{-2}e\wedge e)
   +D_\omega e\wedge\delta\omega\bigr).
\]
The trace pairing is nondegenerate, so the independent variations give both equations.

The chosen matrices satisfy
$\operatorname{Tr}(J_0[J_1,J_2])=1/2$.
Antisymmetrizing the three factors gives
$\operatorname{Tr}(e\wedge e\wedge e)=3\mathrm{vol}_e/2$.
For the curvature term, write
$F_\omega=\tfrac12F^a{}_{bc}J_a\theta^b\wedge\theta^c$.
If $\epsilon^{012}=1$ denotes the alternating symbol, its coefficient is
\[
 \operatorname{Tr}(e\wedge F_\omega)
   =\tfrac14\eta_{ad}F^d{}_{bc}\epsilon^{abc}\,\mathrm{vol}_e,
 \qquad \eta=\operatorname{diag}(-1,1,1).
\]
The curvature acts on the frame by $[F_\omega,J_b]$.
Contracting its matrix coefficients with $\eta^{ab}$ gives
$\mathrm{Scal}=\eta_{ad}F^d{}_{bc}\epsilon^{abc}$.
This last equality can also be checked directly from
$[J_0,J_1]=J_2$, $[J_1,J_2]=-J_0$, and $[J_2,J_0]=J_1$.
The two trace identities give $I_{\mathrm{EH}}$.

For any matrix-valued one-form $v$, expand the two cubic expressions and use the graded cyclic trace.
The result is
\[
 \mathrm{CS}(\omega+v)-\mathrm{CS}(\omega-v)
 =4\operatorname{Tr}(v\wedge F_\omega+\tfrac13v^3)
     -2d\operatorname{Tr}(\omega\wedge v).
\]
For the derivative terms this uses
$d\operatorname{Tr}(\omega\wedge v)
=\operatorname{Tr}(d\omega\wedge v-\omega\wedge dv)$.
Substitute $v=e/\ell$, multiply by $k/(4\pi)$, and use Stokes' formula.
This proves \eqref{eq:pg-action-comparison} with its stated boundary term.
\end{proof}

The normalization depends on the trace and orientation just fixed.
Reversing the orientation changes the sign of every action integral.
Boundary variations that do not vanish require a boundary action and boundary conditions beyond the interior variational problem above.

\section{Boundary transformations and the Virasoro bracket}

The projective transformation law has an infinitesimal form.
For $f(x)=x+\varepsilon\epsilon(x)$, where $\varepsilon^2=0$, equation
\eqref{eq:pg-projective-law} gives
\begin{equation}\label{eq:pg-variation}
 \delta_\epsilon p=\epsilon p'+2\epsilon'p-\tfrac12\epsilon'''.
\end{equation}
The same formula follows from the boundary connection itself.
Let $a=(L_1-p(x)L_{-1})dx$ and let $\lambda$ be a traceless matrix function.
An infinitesimal gauge transformation is $\delta a=d\lambda+[a,\lambda]$.
To retain the coefficient of $L_1$ and keep the $L_0$ coefficient zero, write
\[
 \lambda_\epsilon
 =\epsilon L_1-\epsilon'L_0+(\tfrac12\epsilon''-p\epsilon)L_{-1}.
\]
The $L_1$ and $L_0$ components of $d\lambda+[a,\lambda]$ then vanish.
Its $L_{-1}$ component is $\epsilon'''/2-\epsilon p'-2\epsilon'p$.
This proves \eqref{eq:pg-variation} by an actual connection transformation.

The coefficient of its third derivative can be normalized from the action.
Consider one Chern--Simons sector on a spatial surface $\Sigma$ with an oriented boundary circle of coordinate $x\in\mathbb R/(2\pi\mathbb Z)$.
The boundary orientation is the one in Stokes' formula.
The variation identity
\[
 \delta\mathrm{CS}(A)=2\operatorname{Tr}(\delta A\wedge F_A)
                      -d\operatorname{Tr}(A\wedge\delta A)
\]
gives the one-form on the space of spatial connections
$\Theta_\Sigma(\delta A)=-k\int_\Sigma\operatorname{Tr}(A\wedge\delta A)/(4\pi)$.
Use the convention $\Omega_\Sigma=-d\Theta_\Sigma$ on that space.
Its alternating pairing is
\[
 \Omega_\Sigma(\delta_1A,\delta_2A)
   =\frac{k}{2\pi}\int_\Sigma\operatorname{Tr}(\delta_1A\wedge\delta_2A).
\]
Restrict to flat spatial connections and their tangent variations, so $D_A\delta A=0$.
Integration by parts gives
\begin{equation}\label{eq:pg-boundary-hamiltonian}
 \Omega_\Sigma(D_A\lambda,\delta A)
 =\frac{k}{2\pi}\oint\operatorname{Tr}(\lambda\,\delta a).
\end{equation}
Gauge transformations vanishing on the boundary are null directions of this pairing.
Taking their quotient is the boundary phase space used here.
These statements concern any region of flat-connection space on which the indicated boundary variations extend.
A function $Q$ is a Hamiltonian for a variation $X$ when $\Omega_\Sigma(X,\delta A)=\delta Q$ for every allowed variation $\delta A$.
This condition determines a generator up to null directions of the pairing.

\begin{proposition}\label{prop:pg-virasoro}
On the projective boundary form, the Hamiltonian for a fixed smooth periodic function $\epsilon$ is
\[
 Q_\epsilon=\frac{k}{2\pi}\oint\epsilon p\,dx.
\]
With $\{Q_\epsilon,Q_\zeta\}=\delta_\zeta Q_\epsilon$, its bracket is
\begin{equation}\label{eq:pg-virasoro-bracket}
 \{Q_\epsilon,Q_\zeta\}
 =Q_{\epsilon\zeta'-\zeta\epsilon'}
   -\frac{k}{4\pi}\oint\epsilon\zeta'''\,dx.
\end{equation}
The projective density $\mathcal L=kp$ transforms with central coefficient $c/12$, where
$c=6k=3\ell/(2G)$.
\end{proposition}
\begin{proof}
Since $\operatorname{Tr}(L_1L_{-1})=-1$ and $\delta a=-\delta p\,L_{-1}dx$,
the right side of \eqref{eq:pg-boundary-hamiltonian} is
$k\oint\epsilon\delta p\,dx/(2\pi)$.
This is $\delta Q_\epsilon$.
Although $\lambda_\epsilon$ contains $p$, the Hamiltonian parameter $\epsilon$ is held fixed.
The displayed identity therefore establishes integrability of the charge variation.

Insert \eqref{eq:pg-variation} in $\delta_\zeta Q_\epsilon$ and integrate the $p'$ term by parts around the circle.
The coefficient of $p$ is $\epsilon\zeta'-\epsilon'\zeta$, giving the first term of \eqref{eq:pg-virasoro-bracket}.
The remaining term is its displayed third-derivative integral.
Three integrations by parts make it antisymmetric.
The cocycle identity follows by inserting
$[\epsilon,\zeta]=\epsilon\zeta'-\zeta\epsilon'$ and integrating once more; terms with total derivative order four cancel cyclically.
Equivalently, on the Fourier functions $e^{inx}$ the central term is proportional to
$n^3\delta_{m+n,0}$, whose cyclic identity reduces to
$(m-n)(m+n)^3+(n-l)(n+l)^3+(l-m)(l+m)^3=0$ when $m+n+l=0$.
For smooth periodic functions, repeated integration by parts makes the Fourier coefficients decay faster than every inverse power.
Their differentiated Fourier series converge uniformly, so the identity extends from trigonometric polynomials to these functions.
Thus the bracket is the central extension of circle vector fields called the Virasoro bracket.
Finally, multiplying \eqref{eq:pg-variation} by $k$ gives the coefficient $-k/2=-c/12$.
\end{proof}

For $Q_n=Q_{e^{inx}}$, the same convention reads
\[
 i\{Q_m,Q_n\}=(m-n)Q_{m+n}+\frac{k}{2}m^3\delta_{m+n,0}.
\]
The shift $\widehat Q_n=Q_n+(k/4)\delta_{n,0}$ changes the central term to
$\tfrac{k}{2}(m^3-m)\delta_{m+n,0}$.
This is a classical Poisson algebra; no quantized representation has been assumed.
The critical affine level $-2$ and the gravitational coupling $k=\ell/(4G)>0$ enter different constructions.
The shared projective coefficient alone does not identify them.

\section{Periodic geometry and a regular horizon}

Suppose $p$ and $q$ are smooth $2\pi$-periodic functions.
With $x^+=\varphi+t/\ell$ and $x^-=\varphi-t/\ell$, the identification
$\varphi\sim\varphi+2\pi$ preserves \eqref{eq:pg-metric}.
Periodicity bounds both functions, so one radial lower bound gives a nondegenerate collar for every $t\in\mathbb R$.
The horizontal equations determine its two holonomies around the circle.
A holonomy is the linear map obtained by parallel transport around the loop.
Changing a single-valued frame conjugates that map.
Radial gauge changes therefore preserve its conjugacy class.
This gives a global cylindrical collar with specified monodromy, rather than a choice of a complete interior filling.

Constant profiles expose what can happen at the excluded radial determinant.
Set $p=a^2$ and $q=b^2$, with $a,b>0$.
Define an areal radius $R$ by the proper circumference $2\pi R$ of a circle at fixed $t,r$:
\[
 R^2=\ell^2(s^2+p+q+pq\,s^{-2}).
\]
On the exterior $s^4>pq$, differentiation gives $dR/dr=\ell^2(s^2-pq\,s^{-2})/R>0$.
Introduce
\[
 \begin{gathered}
 R_+=\ell(a+b),\qquad R_-=\ell|a-b|,\\
 N^2(R)=\frac{(R^2-R_+^2)(R^2-R_-^2)}{\ell^2R^2},\qquad
 W(R)=\frac{\ell(p-q)}{R^2}.
 \end{gathered}
\]
Substitution in \eqref{eq:pg-metric} gives
\begin{equation}\label{eq:pg-areal}
 ds_g^2=-N^2dt^2+\frac{dR^2}{N^2}+R^2(d\varphi+Wdt)^2.
\end{equation}
For example, the radial identity is
$N^2=\ell^2(s^2-pq\,s^{-2})^2/R^2$.
The coefficient of $dt\,d\varphi$ is $2\ell(p-q)$ in both expressions.
The remaining coefficients give
$N^2=R^2/\ell^2-2(p+q)+\ell^2(p-q)^2/R^2$, whose factorization is the one displayed above.

A Killing vector is a vector field $K$ with $\mathcal L_Kg=0$.
A Killing horizon is a null hypersurface to which a nonzero Killing vector is normal.
On it, the coefficient in $\nabla_KK=\varkappa K$ is the surface gravity.
Rescaling $K$ by a constant rescales $\varkappa$ by that constant.
Fix the normalization $K=\partial_t-W(R_+)\partial_\varphi$ in the exterior.

\begin{proposition}\label{prop:pg-horizon}
The metric \eqref{eq:pg-areal} extends smoothly and nondegenerately through $R=R_+$.
That hypersurface is a Killing horizon with circumference $2\pi R_+$ and surface gravity
\[
 \varkappa=\frac{R_+^2-R_-^2}{\ell^2R_+}
          =\frac{4ab}{\ell(a+b)}.
\]
\end{proposition}
\begin{proof}
On the exterior, choose antiderivatives and define coordinates by
\[
 dv=dt+N^{-2}dR,\qquad d\vartheta=d\varphi-WN^{-2}dR.
\]
They are locally valid because $N$ and $W$ depend only on $R$.
The angle $\vartheta$ remains periodic with period $2\pi$.
Equation \eqref{eq:pg-areal} becomes
\begin{equation}\label{eq:pg-ingoing}
 ds_g^2=-N^2dv^2+2\,dv\,dR+R^2(d\vartheta+Wdv)^2.
\end{equation}
Its determinant is $-R^2$.
All coefficients are smooth near the positive number $R_+$, including across its zero of $N^2$.
Thus \eqref{eq:pg-ingoing} defines the claimed extension.
Its curvature equation extends by continuity from the exterior.

The chosen vector extends as $K=\partial_v-W(R_+)\partial_\vartheta$.
It preserves the metric because its coefficients are constant in $v,\vartheta$.
On $R=R_+$ it is null, and its metric dual is $dR$.
It is therefore normal to that null hypersurface.
Thus this hypersurface is a Killing horizon.
For a Killing vector, differentiating $g(K,K)$ gives
$\nabla_KK=-\tfrac12\operatorname{grad}(g(K,K))$.
Here $g(K,K)=-N^2+R^2(W-W(R_+))^2$.
At the horizon, the squared term and its derivative vanish.
Since $\operatorname{grad}R=K$ there, the equation is
$\nabla_KK=\tfrac12(N^2)'(R_+)K$.
This coefficient is the surface gravity for the chosen normalization of $K$.
Differentiating the factored $N^2$ gives the displayed value.
The angular line element at fixed $v,R_+$ is $R_+^2d\vartheta^2$, giving the circumference.
\end{proof}

The zero of \eqref{eq:pg-triad} in these radial coordinates has therefore admitted a regular metric extension.
For $a=0$ or $b=0$ with $a+b>0$, the same ingoing formula remains smooth at the coincident roots and has zero surface gravity.
If $a=b=0$, the candidate horizon has $R=0$, where its determinant also vanishes; the preceding positive-radius extension does not apply.

\section{Gauge transformations and invertible triads}

An invertible triad also controls the relation between infinitesimal gauge transformations and changes of coordinates.
Let $\epsilon^+,\epsilon^-$ be their two matrix parameters.
There is a unique vector field $\xi$ such that
\[
 e(\xi)=\tfrac\ell2(\epsilon^+-\epsilon^-).
\]
Set $\lambda=(\epsilon^++\epsilon^-)/2-\omega(\xi)$.
Then $\epsilon^\pm=A^\pm(\xi)+\lambda$.
The formula
\[
 \mathcal L_\xi A=\iota_\xi F_A+D_A(\iota_\xi A)
\]
follows from the product rule and the definition of the Lie derivative of a one-form.
For flat $A^\pm$, it gives
\[
 \delta e=\mathcal L_\xi e+[e,\lambda],\qquad
 \delta\omega=\mathcal L_\xi\omega+D_\omega\lambda.
\]
The commutator term is a Lorentz rotation and leaves $g$ invariant.
Hence the metric changes by its Lie derivative along $\xi$.
This is a local infinitesimal comparison on the invertible-triad region.

The larger space of flat connections also contains $A^+=A^-=0$, where $e=0$.
More generally, on a contractible neighbourhood each flat connection can be separately trivialized by solving its horizontal equation.
Their independent gauge transformations can therefore take a nondegenerate pair to a degenerate one.
The inverse triad used to define $\xi$ ceases to exist at that boundary.
This supplies an explicit obstruction to identifying the unrestricted gauge quotient with the space of nondegenerate gravitational metrics.

The maps constructed here retain their additional data:
\[
 \begin{gathered}
 \chi\longmapsto \partial_z^2-\mathcal U,\\
 (p,q,\ell)\longmapsto(A^+,A^-)
       \xrightarrow{\det e\ne0}(g,\nabla).
 \end{gathered}
\]
The first arrow selects a projective equation from a critical character.
The second supplies two real boundary profiles and a radial construction.
The last uses the trace form and an invertible triad.
For periodic profiles these data give a cylindrical collar; particular constant profiles admit the horizon extension just derived.
An interior filling, allowed boundary variations, global gauge identifications, and a quantized state space require further mathematical data and constructions.
They are not determined by a single quadratic central character.
